% RESEARCH_DRAFT_SHA256: 7f5efa776cdc6d35256ce0def31c8f14f5a387c2e36aa7ef129f2d231cb98940
% Layout and source-reference wrappers expanded for the web renderer.
\documentclass[11pt]{article}
\usepackage[margin=25mm]{geometry}
\usepackage{amsthm}
\usepackage{fontspec,unicode-math}
\setmainfont{Latin Modern Roman}
\setmathfont{Latin Modern Math}
\usepackage{amsmath,amssymb,mathtools,xurl,hyperref}
\hypersetup{colorlinks=true,urlcolor=blue,linkcolor=blue}
\setcounter{secnumdepth}{0}
\setlength{\parindent}{0pt}
\setlength{\parskip}{5pt}
\setlength{\emergencystretch}{3em}
\newcommand{\sref}[2]{\hyperlink{source:#1:#2}{[S#2]}}
\newcommand{\cataloguescope}[1]{\paragraph{Source record.}#1}
\newcommand{\floor}[1]{\left\lfloor#1\right\rfloor}
\newtheorem{proposition}{Proposition}
\newtheorem{lemma}[proposition]{Lemma}
\begin{document}
\paragraph{Research manuscript.} OpenAI \texttt{gpt-6-astra}, reasoning \texttt{ultra}; 11 October 2026.
\paragraph{Preserved source body.} The complete catalogue mathematical text and sources follow with layout and source links expanded, including source qualifications and its historical preparation notes. New research begins at the section headed Research attempt.

\section{Certified open-world routing for a database of databases}\label{statement}

\subsection{Definitions and mathematical statement}

\paragraph{Shared notation.}Standard number systems are $\mathbb N,\mathbb Z,\mathbb Q,\mathbb R,\mathbb C$. Specialized notation and index ranges follow the cited source and the definitions below.


A database catalogue records metadata about databases, such as their subject matter, identifiers, access locations, and coverage. A database of databases indexes those catalogues or resources. Discovery retrieves resources relevant to a user's information need, whereas querying a selected database retrieves its underlying data.
\paragraph{Record scope and qualifications.}
The catalogue title is retained as metadata. Named certificates, contracts, bounds or frameworks in that title are not asserted as results of the original source extract.
No resource universe, metadata schema, query language, ranking objective, or completeness guarantee is given. This is an infrastructure project.

Catalogue screening limitation: Workshop planning or software infrastructure agenda; the supplied passage does not identify a determinate mathematical target.

The source snapshot labels this record partially solved; this is catalogue metadata rather than independent verification of a mathematical result or the current completion of a project.
\paragraph{Statement recorded in the supplied source.}
This problem asks generally how to build a database of databases to help users navigate to the resource where the information they wish to obtain is.
The concrete goal of this objective is to create such a database or improve the one that exists and outline how it could be shaped in the future in order to serve the community even better. \href{http://aimpl.org/cyberinfrastructure/1/}{[S2]}
\paragraph{Source status.}
The supplied catalogue labels this source record partially solved; that metadata was checked against the supplied snapshot on 2026-10-11. The label is preserved for provenance and is not an independent verification of present open status or mathematical correctness.
\subsection{Short English statement}
Build or improve a catalogue that directs users to appropriate mathematical databases.
\subsection{Sources}
\begin{itemize}
\item[S1] Ulam AI and the UnsolvedMath Contributors, UnsolvedMath: A Curated Collection of Open Mathematics Problems, catalogue entry. Catalogue curation and metadata: CC BY 4.0. The source extract is retained; definitions, English summaries and qualifications are editorial additions.
\url{https://huggingface.co/datasets/ulamai/UnsolvedMath}.
\item[S2] AIM Workshop Problem Lists, source linked by the catalogue.
\url{http://aimpl.org/cyberinfrastructure/1/}.
\end{itemize}
\paragraph{Preparation and review.}
The original source and UnsolvedMath contributors retain attribution for the source extract. Codex (GPT-6; exact serving revision not exposed) assisted with definitions, English summaries, scope notes and formatting. No new human mathematical authorship or independently verified solution is asserted. Catalogue identity checked against the complete supplied registry and current repository. Distinct upstream records, including repeated questions and broad or damaged extracts, are retained for complete coverage of the source CS category. A repeated record is not a new independent conjecture. Complete mathematical-equivalence and current-literature review has not been performed. Preparation and final review follow contribution prompts C and D, with source gaps disclosed. Structural and rendering checks do not establish a new theorem or open status.
\label{end:statement}

\section{Research attempt}
\subsection{Literature, source context and scope}
The complete catalogue body above retains the original agenda, source status, and historical preparation. The AIMPL problem page could not be retrieved. The supplied complete source transcription was read, and the accessible primary AIM workshop announcement was checked for context. The announcement is not a substitute for verification of the inaccessible problem page. The catalogue's earlier proposed frameworks and research summaries are context, not established solutions credited to this attempt. The following is an explicitly limited investigation of one route; it neither substitutes a formalized toy model for the entire source agenda nor claims empirical or deployed-software results. Elementary auxiliary facts are not new original-target progress.

\subsection{Testing complete routing from incomplete resource metadata}
The requested database of databases is an infrastructure project. We examine a specific proposed shortcut: decide conclusively which resources contain an answer using only topic metadata. In a finite model a directory observes a resource's name and topic label but not its full contents; a query asks whether record 17 occurs. Two possible underlying worlds have identical visible metadata and resource identities. In world $W_0$ the resource contains records $\{2,5\}$, and in $W_1$ it contains $\{2,5,17\}$.
\begin{proposition}
No rule depending only on that visible directory can always return the correct binary answer to this query. More generally, a negative answer is certified by visible information only if every world consistent with that information lacks an answer.
\end{proposition}
\begin{proof}
The rule receives the same input in both worlds and hence has the same output. Correct answers are respectively no and yes, so a deterministic output fails in one world. A randomized rule with yes-probability $p$ has error $p$ in $W_0$ and $1-p$ in $W_1$, at least one at least $1/2$. For the second statement, if a consistent world contains an answer, declaring no cannot be correct in all consistent worlds. Conversely if all consistent worlds lack one, no is correct in every such world by definition.
\end{proof}
A sound limited router can therefore use three results: verified hit if all admissible worlds contain a hit, verified miss if all lack one, and unknown otherwise. This is a semantic specification, not an efficient algorithm for an unlimited collection of worlds. In the two-world example it returns unknown; inspecting a contents certificate eliminates one world and determines the answer. An open collection may include unlisted resources, so certification of one resource's miss cannot automatically certify global absence.
\subsection{Verification and remaining gaps}
The two resource sets were queried directly, giving opposite answers under identical topic metadata. No live registry was created or improved, and no metadata coverage, certificate format, update policy, or complete search method was implemented. The indistinguishability argument blocks one overstrong routing guarantee; it does not answer the entire infrastructure agenda.

\section{References and source attribution}
Primary source: \url{http://aimpl.org/cyberinfrastructure/1/}. Attempted primary AIMPL page retrieval; failed with timeout or gateway error. The supplied full source transcription was read; the inaccessible live page was not independently inspected.

\url{https://aimath.org/pastworkshops/cyberinfrastructure.html}. Read the primary December 2023 workshop announcement; it describes the open-source mathematical infrastructure agenda.

All retrievals were checked on 11 October 2026. Original contributor and catalogue attribution remains in the preserved Sources section. The elementary arguments below are self-contained; their exposition is not a claim of first discovery.


\paragraph{Integration review.} Codex (GPT-6; exact serving revision not exposed) reviewed the written arguments and expanded layout and source-reference wrappers. Repository metadata and the production web renderer were checked using Python 3.13 and JavaScript with MathJax 3.2.2. This model-assisted review is not independent human proof verification; no PDF compilation or formal proof checking is asserted.

\section{Completion Estimate}
\noindent\textbf{COMPLETION: 0\%}
No substantiated new progress toward the entire original target is claimed. The worked results and failed route are delimited above; a known result, illustrative model, or auxiliary obstruction is not counted as original-target completion.

\section{Final status and attribution}
\textbf{Attempt status: unresolved.} This describes this research attempt, not a new ruling on the source catalogue's status or on external literature claims.
The mathematics and this self-audit were generated by OpenAI \texttt{gpt-6-astra}, reasoning \texttt{ultra}, on 11 October 2026. Codex (GPT-6; exact serving revision not exposed) orchestrated the batch and contributed formatting/repository review. No human mathematical authorship, independent human proof checking, or formal verification is claimed. The historical preparation paragraphs in the reproduced catalogue body concern that body only.

Contribution prompt A was applied to the complete definition and supplied source record, followed by the adversarial checks of prompt B. Local reproduction files are \texttt{batch013/checks.py}, \texttt{computations.json}, and the per-ID audit JSON. They were not public links at preparation time. Numerical checks are finite and do not prove unrestricted results. No LaTeX engine was found locally; compilation was not performed by this research agent.

\end{document}
